\begin{proposition}[P4.10]\hypertarget{tag-02WH}{}\label{night:P4.10}\label{stmt:P4.10}
\StatementRecord{night:P4.10}{proved}{02WH}{the definitions and constructions stated in the proof}

There is a smooth projective threefold $Z$ with $b_1(Z)=0$ and a
connected torus fibration $Z\to\Pp^1$ whose reduced special fibres
are irreducible normal-crossings surfaces, but $c_3[Z]=24$.
Consequently the hypothesis $b_2=0$ cannot be omitted, even in the
intersection of the two classes of fibres in the question.
\end{proposition}
\begin{proof}

\begin{enumerate}[label={$\langle1\rangle$\arabic*.},leftmargin=*]
\item On the two affine charts of $\Pp^1$, with $s=t^{-1}$, we take
\[
 y^2z=x^3-3xz^2+t^6z^3,\qquad
 y_\infty^2z=x_\infty^3-3s^4x_\infty z^2+z^3,
\]
glued by $x_\infty=s^2x$ and $y_\infty=s^3y$. These are the
Weierstrass cubics in $\Pp(\OO\oplus\OO(2)\oplus\OO(3))$ for
the line bundle $\OO_{\Pp^1}(1)$. They define an elliptic surface
$\pi:S\to\Pp^1$ with zero section $[0:1:0]$. The discriminant
is a nonzero constant times $t^{12}-4$, or $1-4s^{12}$ on the
other chart. There are twelve simple zeros and a smooth fibre at
infinity. At a discriminant zero the cubic has a double root
$x=1$ or $x=-1$ and a distinct third root, so the fibre has an
ordinary node. Since $6t^5\ne0$ there, the total surface is smooth.
At the zero section the derivative in $z$ is nonzero. The map near
each node is therefore $xy=t$ after holomorphic coordinate changes
and translation of the base.

\item The relative dualizing differential $dx/y$ has transitions
$\pi^*\OO(1)$ and generates the dualizing sheaf also at the nodes.
Thus $\omega_S=\pi^*\OO(-1)$. Relative duality gives
$R^1\pi_*\OO_S=\OO(-1)$, and connected fibres give
$\pi_*\OO_S=\OO$. Leray yields $h^1(\OO_S)=0$, while all
positive plurigenera vanish. The classification \cite[Chapter~VI, Theorem~1.1, p.~243]{BHPV2004E} makes $S$
rational; in the ruled case the base genus equals $h^1(\OO_S)=0$.
In particular $\pi_1(S)=1$, since the minimal rational models
$\Pp^2$ and the Hirzebruch surfaces are simply connected and a
point blowup preserves the fundamental group.

\item The projective fibre product $W=S\times_{\Pp^1}S$ has exactly
twelve singularities, at the pair of nodes in each nodal fibre.
Their equations are $xy=uv$, with base function $xy=uv$. We blow
up the coherent ideal of the diagonal $\Delta_S\subset W$. This
divisor is Cartier away from the twelve points, so the blowup is
the identity there. Put $r=x-u$, $h=y-v$ at a singular point.
The equation becomes $ry+uh=0$ and the diagonal ideal is $(r,h)$.
The blowup has charts
\label{layout:II:1}
\[
\begin{gathered}
h=ar,\quad y=-au,\quad x=u+r,\quad f=-aux;
 \qquad
 \\
r=bh,\quad u=-by,\quad x=-bv,\quad f=-byv.
\end{gathered}
\]
The independent coordinates are $(a,u,x)$ and $(b,y,v)$.
Both charts are smooth, with $b=a^{-1}$ on the overlap, and the
exceptional fibre is $\Pp^1$. Thus the resulting $Z$ is smooth
and projective. The displayed equations make every special fibre
reduced and normal crossings. Its support remains irreducible:
the original product of two irreducible nodal cubics is irreducible,
and the exceptional locus has dimension one, so introduces no
surface component. Its normalization is rational, by the
birational map from $\Pp^1\times\Pp^1$.

\item A nodal cubic has Euler number $2-1=1$, so its self-product has
Euler number one. The resolution replaces one point by $\Pp^1$;
each resolved special support therefore has Euler number two.
All other fibres are $E_t\times E_t$ with Euler number zero.
Hence $e(Z)=12\cdot2=24=c_3[Z]$.

\item For completeness we compute $\pi_1(Z)$. Every loop can be moved
off the fibre divisors, so the fundamental group of their complement
surjects onto $\pi_1(Z)$. The complement is a bundle with fibre
$E\times E$ and a section over a twelve-punctured sphere. The
section extends over $\Pp^1$, killing the base generators. The
two maps
\[
 S\longrightarrow Z:\quad p\mapsto(p,O(\pi(p))),\qquad
 p\mapsto(O(\pi(p)),p)
\]
are defined because their images miss every conifold point, and
extend the two factor inclusions of $E$. Since $S$ is simply
connected they kill both fibre fundamental groups. Thus
$\pi_1(Z)=1$ and $b_1(Z)=0$. A projective Kähler class is nonzero
in $H^2(Z;\R)$, because its cube has positive integral. This
identifies the omitted hypothesis: $b_2(Z)>0$.

\end{enumerate}
\end{proof}